% Resumen (espanol).
\chapter*{Resumen}
\addcontentsline{toc}{chapter}{Resumen}
\markboth{Resumen}{Resumen}

Esta tesis presenta un modelo reproducible de crecimiento urbano para la Zona
Metropolitana de Querétaro, construido a partir de una serie anual de imágenes
de 37 años (1984--2020) y del acoplamiento entre Pesos de Evidencia (WoE) y
autómatas celulares (AC). Las capturas RGB obtenidas con la herramienta de
imagen histórica de Google Earth se transformaron en mapas binarios
urbano/no-urbano mediante un \textit{pipeline} de clasificación basado en
patrones binarios locales, agrupamiento K-Means y un clasificador SVM lineal.
Sobre esos mapas, los Pesos de Evidencia cuantifican la influencia espacial de
un conjunto reducido de variables auditables y alimentan, junto con la vecindad
de Moore, una regla de transición probabilística que el autómata celular aplica
celda a celda.

El modelo se calibró mediante búsqueda sistemática sobre el umbral de
transición, el peso de la vecindad y la ponderación por \textit{Information
Value}, y se evaluó con un protocolo de validación en cinco ventanas
quinquenales independientes (2011--2016 a 2015--2020). El desempeño promedio
alcanza un \textit{Figure of Merit} de $0{,}317$, un coeficiente Kappa de
$0{,}445$, una exactitud de $0{,}722$ y un \textit{IoU} de $0{,}633$, con la
descomposición de error de Pontius como marco de interpretación. La
contribución central no es el acoplamiento WoE-AC en sí, sino un protocolo de
validación multitemporal reproducible, aplicado a un contexto mexicano poco
documentado, con una implementación de caja blanca que corre íntegramente en
CPU, sin GPU ni software propietario, a partir de datos de acceso libre.

\vspace{0.6cm}
\noindent\textbf{Palabras clave:} crecimiento urbano; autómatas celulares;
Pesos de Evidencia; validación multitemporal; \textit{Figure of Merit};
Zona Metropolitana de Querétaro.
\cleardoublepage
